\documentclass[11pt]{article}
\usepackage[margin=1in]{geometry}
\usepackage[T1]{fontenc}
\usepackage[utf8]{inputenc}
\usepackage{lmodern}
\usepackage{setspace}
\usepackage{parskip}
\setlength{\parindent}{0pt}
\setstretch{1.05}

\begin{document}

Decory Edwards \\
Baltimore, MD \\
\today

\vspace{1em}

Search Committee \\
Federal Reserve Bank of Richmond \\
Research Department

\vspace{1em}

Dear Members of the Search Committee,

I am writing to apply for the Regional Economist position in the Research Department at the Federal Reserve Bank of Richmond. I am a PhD candidate in Economics at Johns Hopkins University, and my training has prepared me for a role that combines empirical research, clear communication, and public-facing economic analysis. The opportunity to contribute to policy-relevant regional research in Baltimore is especially appealing because it would allow me to bring together my quantitative background, my interest in economic storytelling, and my connection to the local community.

My dissertation draft studies how heterogeneity in financial environments shapes household wealth accumulation and realized investment performance. In the first essay, I develop a heterogeneous-agent model with dispersion in household asset returns and show that return heterogeneity substantially improves the model's ability to match empirical wealth moments, while also comparing the distributional and welfare implications of wealth and capital income taxation. In the second essay, I use the RAND Longitudinal HRS to measure returns to net wealth and document a hump-shaped relationship between generalized trust and realized investment performance after accounting for demographic characteristics and portfolio composition. Together, these projects reflect my interest in combining theory, data analysis, and careful empirical interpretation to understand meaningful economic variation across households and environments.

I am also strongly interested in regional research. I've been a student and consumer in the Baltimore community for six years now. As a trained economist, I believe I can use my interest in analyzing trends in data and testing economic theories to benefit the Baltimore, DC, and West Virginia areas. I am especially drawn to a role that would let me study regional conditions closely, communicate findings to a range of audiences, and contribute to the Bank's outreach and survey work.

This position fits naturally into my career path. I want to continue building the human capital I have developed in understanding and crafting economic arguments through employment as a Regional Economist. The Richmond Fed role would let me deepen my empirical skills, strengthen my ability to translate economic developments for public and policy audiences, and continue growing as an economist whose work connects rigorous evidence to real regional and national questions.

Thank you for your time and consideration. I would be grateful for the opportunity to discuss my application further.

Sincerely,

\vspace{1em}

Decory Edwards

\end{document}
